%% ma: L10-B13-0025
%% lop: 10
%% chuong: 5
%% bai: 13
%% dang: 10.13.3
%% loai: TN4
%% muc: TH
%% dap_an: "D"
%% nguon: "chuyen-de-cac-so-dac-trung-cua-mau-so-lieu-khong-ghep-nhom-toan-10-chuong-trinh-moi.pdf (GV. Phan Nhật Linh) (P02), Bài 02 Dạng 1 Phần I Câu 15"
%% kien_thuc: "Lớp 10 Bài 13 (trung vị)"
%% trang_thai: chinh-thuc
%% ghi_chu: "Tự duyệt theo yêu cầu giáo viên 22/10/2026."
\begin{ex}
Điều tra số học sinh của $30$ lớp học, ta được bảng số liệu (đã sắp xếp) như sau:
\begin{center}
\begin{tabular}{|c|c|c|c|c|c|c|c|c|c|c|c|c|c|c|}\hline
$35$ & $39$ & $39$ & $40$ & $40$ & $41$ & $41$ & $41$ & $41$ & $44$ & $44$ & $45$ & $45$ & $45$ & $46$\\\hline
$48$ & $48$ & $48$ & $48$ & $49$ & $49$ & $49$ & $49$ & $49$ & $49$ & $50$ & $50$ & $50$ & $50$ & $51$\\\hline
\end{tabular}
\end{center}
Trung vị của mẫu số liệu trên là
\choice{$46$}{$48$}{$45$}{\True $47$}
\loigiai{$n=30$; giá trị thứ $15$ và thứ $16$ lần lượt là $46$ và $48$ nên $M_e=\dfrac{46+48}2=47$. Chọn D.}
\end{ex}
